% ============================================================================
\section{Discussion and Conclusion}
\label{sec:conclusion}
% ============================================================================

The local analysis distinguishes the expected effect of a shared update
from the uncertainty in that effect. When a task's progress margin is
nonpositive, changing sample counts at fixed weights cannot make its
expected first-order progress positive under the frozen-preconditioner
model. For positive margins, additional samples can reduce uncertainty.
GPAS allocates samples to reduce total update variance, which bounds
directional variance conservatively. Directional variance determines a more specific
bound for each task; the two allocation criteria need not agree.

The proposed study tests these quantities against updates from identical
checkpoints and against complete training runs. Separate-distillation
references assess individual transfer, while fixed-state loss, fresh-policy
loss, and capability evaluation measure different aspects of training.
\missing{Report which diagnostics predicted observed failures, when GPAS
improved capability or compute efficiency, and which failures remained.}

These local conditions do not establish joint capability attainability or
global convergence. Successful training may temporarily increase a task
loss, and a local conflict does not prove a capacity limit. The empirical
scope is one student size, four domains, and one hardware layout; the
repeated core comparison estimates training variation within that scope.
Finite-sample diagnostic error, changing optimizer states, and changing
rollout distributions must be assessed before using the local bounds to
make training decisions.
